%======================================================================
\section{The point-Pauli game: a polynomial window}\label{sec:point}
%======================================================================

The configuration lamps of \cref{sec:lamps} sit on the $2^N$ configurations of $N$ points and give an almost inverse-loss
demand. Placing Pauli lamps on the points themselves gives a fixed game whose demand lies between two positive powers of $1/\eps$,
strictly between the logarithmic and inverse scales. Whether its demand has a limiting power-law exponent, and what that exponent is,
remain open.

\subsection{The source and the game}\label{sec:point-source}

The \emph{point source} has generators $a,b,s,x,z,J$ and the $22$ relators
\begin{gather*}
s^2,\quad(sas a^{-1})^3,\quad[s,a^2sa^{-2}],\quad aba^{-1}b^{-1}s^{-1}\ (\mathrm{H4}),\quad asa^{-1}bs^{-1}b^{-1},\\
J^2,\ x^2,\ z^2,\quad[J,a],[J,b],[J,s],[J,x],[J,z],\quad[x,z]J^{-1},\\
[x,sa],\ [x,sb],\ [z,sa],\ [z,sb],\qquad[u,ava^{-1}]\quad(u,v\in\{x,z\}),
\end{gather*}
with $[u,v]=uvu^{-1}v^{-1}$ and marked involution $J$. The first five are Johnson's presentation of Houghton's group
$H_3$~\cite{Joh99},~\cite[Theorem~2.8]{Lee12}, in the words of~\cite[(1)]{DM-H3}; $s$ is the transposition of the first two points of
the first ray, $k=sa$ and $l=sb$ generate the stabilizer of the first point, and the last eight relators place a Pauli pair
$x,z$ at that point, commuting with its stabilizer, with all sites' pairs commuting and every onsite commutator equal to the
central $J$. Let $G_{\mathrm{pt}}$ be the game of the universal compiler (\cref{thm:universal}) for this presentation, with the
compiler's uniform-context, uniform-incidence question law. It has $6960$ variables and $5713$ contexts, of arity at most three,
and upper amplification $12$. The exported game is archived in \texttt{verify/artifacts}, where an independent checker verifies,
exactly and against these $22$ relators, every variable square, context commutator and forbidden assignment and the six
recovery maps.

\subsection{Fillings on the actual state}\label{sec:point-filling}

The lower bound needs the cost, on the actual state, of long null words with large fillings. \Cref{lem:conjugators} and
\cref{lem:calculus}(d) charge each of the $A$ cells of a filling its own conjugator, of length up to $n+K(A+1)$ for a word of
length $n$, and so give the quadratic bound $A(1+2n+2K(A+1))e$. The conjugators can be shared. Cut a van Kampen diagram into a
tree of faces; at each face, substitute at once the whole boundary of the subdiagram that hangs from each child edge, and pay only
for moving the bounded remainder of that one face. Every edge of the tree is then charged once, and the charge becomes linear in
the size of the diagram (\cref{thm:linear-filling}).

In this subsection $\varphi:F(S)\to\mathrm U(d)$ is a homomorphism, $\xi=\rho^{1/2}$ for a density matrix $\rho$, and
$\hs{[\varphi(s),\xi]}\le e_c$ for every $s\in S$. Words may contain the letter $1$, with $\varphi(1)=I$ and $1^{-1}=1$, and $|v|$
counts the letters of a word $v$ other than $1$.

\begin{lemma}[Bounded transport]\label{lem:bounded-transport}
For every operator $Y$ and every word $v$, $\hs{Y\varphi(v)\xi}\le\hs{Y\xi}+\opn Y\,|v|\,e_c$.
\end{lemma}

\begin{proof}
$Y\varphi(v)\xi=Y\xi\varphi(v)+Y[\varphi(v),\xi]$. The first term has norm $\hs{Y\xi}$ because $\varphi(v)$ is unitary, and the
second at most $\opn Y|v|e_c$ by \cref{lem:calculus}(a).
\end{proof}

The bounded operator $Y$ is moved past $\varphi(v)$ as a whole, so the length of the word that produced $Y$ never enters. With
actual mirrors on a purification $\Psi$ of the original state, \cref{lem:transfer-parties}(b) is the same statement, with $e_c$
replaced by the mirror error of the signed generators.

\begin{lemma}[Substituting whole words into one face]\label{lem:face-substitution}
Let $t_1,\dots,t_m$ be letters of $S^{\pm1}\cup\{1\}$ and $q_1,\dots,q_m$ words, and put $R=t_1\cdots t_m$, $W=q_1\cdots q_m$ and
$b_i=\hs{(\varphi(t_i^{-1}q_i)-I)\xi}$. Then
\[
\hs{(\varphi(W)-I)\xi}\ \le\ \hs{(\varphi(R)-I)\xi}+\sum_{i:\,q_i\ne t_i}\bigl(b_i+2(m-i)e_c\bigr).
\]
\end{lemma}

\begin{proof}
Telescoping,
\[
\varphi(W)-\varphi(R)=\sum_{i=1}^m\varphi(q_1\cdots q_{i-1})\bigl(\varphi(q_i)-\varphi(t_i)\bigr)\varphi(t_{i+1}\cdots t_m),
\]
and the terms with $q_i=t_i$ vanish. The left factors are unitary. The middle factor
$Y_i=\varphi(t_i)\bigl(\varphi(t_i^{-1}q_i)-I\bigr)$ has $\opn{Y_i}\le2$ and $\hs{Y_i\xi}=b_i$, so \cref{lem:bounded-transport}
bounds the $i$-th term, applied to $\xi$, by $b_i+2(m-i)e_c$. Finally
$\hs{(\varphi(W)-I)\xi}\le\hs{(\varphi(R)-I)\xi}+\hs{(\varphi(W)-\varphi(R))\xi}$.
\end{proof}

The words $q_i$ may be arbitrarily long; only the suffix of the one face $R$ is moved.

A \emph{disk diagram} is a finite connected plane graph whose unbounded face is bounded by a simple cycle, with every edge
oriented and labelled by a letter of $S^{\pm1}\cup\{1\}$; an edge read against its orientation gives the inverse letter. Its
\emph{faces} are the bounded faces of the graph. The boundary walk of a face, or of the whole diagram, read counterclockwise
from one of its edges is a word, its \emph{boundary word} from that edge. An edge on two faces is read in opposite directions by
them, and an outer edge is read in the same direction by the diagram and by its face.

\begin{lemma}[A tree of faces]\label{lem:face-tree}
Let $D$ be a disk diagram with $F\ge1$ faces, each with a boundary walk of at most $K'$ edges, and let $w$ be its boundary word
from an outer edge $g_0$. Suppose that every face $f$ has $\hs{(\varphi(u)-I)\xi}\le r_f$ for its boundary word $u$ from each of
its edges. Then
\[
\hs{(\varphi(w)-I)\xi}\ \le\ \sum_fr_f+2K'(F-1)e_c .
\]
\end{lemma}

\begin{proof}
\emph{A tree of faces.} The outer boundary is a simple cycle $Z$. The path $Z\setminus g_0$ is acyclic, so it extends to a
spanning tree $T$ of the graph, and $g_0\notin T$. The edges outside $T$ are dual to a spanning tree $T^*$ of the dual graph,
whose vertices are the $F$ faces and the unbounded face. Indeed, a cycle of such dual edges is a closed curve that crosses each
corresponding primal edge once and meets the graph nowhere else, so by the Jordan curve theorem it separates the two ends of each
such edge; the path in $T$ between them would have to cross the curve, through an edge outside $T$, which is impossible. So these
dual edges form a forest, and by Euler's formula there are $E-V+1=F$ of them, one fewer than the number of dual vertices; hence
$T^*$ is a spanning tree. The only outer edge outside $T$ is $g_0$, so the unbounded face is a leaf of $T^*$, attached to the face
$f_0$ that contains $g_0$. Deleting it leaves a tree $T^*_0$ on the $F$ faces, whose $F-1$ edges are interior edges outside $T$.
Such an edge lies on two distinct faces, since an edge with one face on both sides is a bridge and every bridge lies in $T$, and
each of the two faces traverses it once.

\emph{Substituting subtrees.} Root $T^*_0$ at $f_0$. For $f\ne f_0$ let $g_f$ be the edge joining $f$ to its parent. Let
$R_f=t_1\cdots t_m$, $m\le K'$, be the boundary word of $f$ from $g_f$, and that of $f_0$ from $g_0$. Define $W_f$, recursively
from the leaves, by replacing in $R_f$ each letter $t_i$ read along the edge $g_c$ of a child $c$ by the word $q_i$ obtained from
$W_c$ by deleting its first letter. That first letter is read along $g_c$ in the direction opposite to $f$, so it is $t_i^{-1}$,
and $t_i^{-1}q_i=W_c$. Put $\beta_f=\hs{(\varphi(W_f)-I)\xi}$. By \cref{lem:face-substitution}, with $m-i<K'$,
\[
\beta_f\ \le\ r_f+\sum_{c\text{ child of }f}\bigl(\beta_c+2K'e_c\bigr),
\]
and induction from the leaves gives $\beta_{f_0}\le\sum_fr_f+2K'(F-1)e_c$, since $T^*_0$ has $F-1$ edges.

\emph{Free equality.} Read each $W_f$ as an edge path. The boundary walk of $f$ from $g_f$ is closed; replacing its traversal of
$g_c$, from $u$ to $v$ say, by the remainder of the closed path $W_c$, which runs from $u$ back to $v$, keeps it a closed path. So
$W_{f_0}$ is a closed edge path beginning with $g_0$. Every edge $g_c$ of $T^*_0$ is traversed once by its parent and once by its
child, at the first letter of $W_c$, and both traversals are deleted; so $W_{f_0}$ uses only edges of $T\cup\{g_0\}$ and traverses
$g_0$ exactly once, in the direction of $w$. In the graph $T\cup\{g_0\}$, whose only cycle is $Z$, a closed edge path from the
initial vertex of $g_0$ is determined up to homotopy by the signed number of its traversals of $g_0$. Hence $W_{f_0}$ and the
boundary walk $w$ are homotopic, they become equal after cancelling backtracking pairs of edges, and their words are freely equal.
So $\varphi(w)=\varphi(W_{f_0})$.
\end{proof}

A van Kampen diagram may be singular, with edges and vertices on no face, and a face may meet itself; the next lemma removes this
at linear cost.

\begin{lemma}[Regular diagrams at linear cost]\label{lem:regular-diagram}
Let $R$ be a set of words of length at most $K\ge1$, and let a freely reduced word $w_0\ne1$ of length $n$ be freely equal to a
product of $A$ conjugates of elements of $R^{\pm1}$. There is a disk diagram whose boundary word from a suitable edge is $w_0$
with letters $1$ inserted, and whose faces are at most $A$ \emph{caps} and further \emph{auxiliary} faces. The boundary word of a cap, from
any of its edges, is a cyclic permutation of an element of $R^{\pm1}$ with letters $1$ inserted, of at most $2K$ letters; that of
an auxiliary face has at most four letters and is freely trivial. There are at most $A+5E$ faces, where $E$ is an integer with
$2E\le KA+n$.
\end{lemma}

\begin{proof}
By van Kampen's lemma~\cite[Chapter~V, Section~1]{LS77}, as in the proof of \cref{lem:conjugators}, $w_0$ has a van Kampen
diagram $D_0$ with at most $A$ faces, the boundary walk of each reading a cyclic permutation of an element of $R^{\pm1}$, and with outer
boundary walk $w_0$ from a base vertex. Let $E$ be its number of edges. Each edge has two sides, and each side lies on exactly one
face boundary walk or on the outer walk, so $2E\le KA+n$.

Thicken the $1$-skeleton of $D_0$ in the plane: a small closed disk around each vertex and a thin band along each edge, meeting
the disks of its two ends in two short arcs. The union $N$ is a connected compact planar surface, and the bounded components of
its complement are the faces of $D_0$, slightly shrunk, one for each. Fill each of them with a disk, the cap of that face; by the
Jordan--Schoenflies theorem, $N$ with its caps is a closed disk, bounded by the outer boundary curve of $N$. The two long sides of
the band of an edge carry the letter of that edge, oriented alongside it, and every short band end, every arc of a vertex disk and
every diagonal added below carries $1$. A band then reads $t\,1\,t^{-1}\,1$, which is freely trivial from every edge. A cap meets,
along its boundary, one long band side for each edge of its face and one vertex arc between consecutive ones, so it reads the
relator of that face with letters $1$ inserted. The outer boundary meets one long band side for each step of the outer walk of
$D_0$, including both sides of an edge that the walk traverses twice, so, read from the band side of the first edge of $w_0$, it is
$w_0$ with letters $1$ inserted. The disk of a vertex of degree $\delta\ge2$ is a polygon with $2\delta$ sides, alternately band ends
and free arcs; divide it by diagonals into $2\delta-2$ triangles. The disk of a vertex of degree one has two sides and is kept.
These faces read words in the letter $1$ only. As the degrees sum to $2E$, there are at most $A+E+4E$ faces.
\end{proof}

\begin{theorem}[A linear filling charge]\label{thm:linear-filling}
Let $R$ be a finite set of words of length at most $K\ge1$, and put $K'=\max\{2K,4\}$. Suppose that $\hs{[\varphi(s),\xi]}\le e_c$
for every $s\in S$ and $\hs{(\varphi(r)-I)\xi}\le e_r$ for every $r\in R$. If a word $w$ of length $n$ has area at most $A$ in $R$,
then
\[
\hs{(\varphi(w)-I)\xi}\ \le\ Ae_r+\bigl[(2K+2K'+5KK')A+5K'n\bigr]e_c .
\]
In particular, in a model at precision $e$ the error is at most $C_K(A+n)e$, with $C_K=1+2K+2K'+5K'(K+1)$.
\end{theorem}

\begin{proof}
Let $w_0$ be the free reduction of $w$, so that $\varphi(w)=\varphi(w_0)$ and $|w_0|\le n$; if $w_0=1$ there is nothing to prove.
Otherwise take the disk diagram of \cref{lem:regular-diagram} and apply \cref{lem:face-tree}, whose faces have at most $K'$
edges. A cap read from any edge is, up to letters $1$, a word $\alpha^{-1}r^{\pm1}\alpha$ with $r\in R$ and $|\alpha|\le K$. Since
$\varphi(r)^{-1}-I=-\varphi(r)^{-1}(\varphi(r)-I)$, \cref{lem:calculus}(c) and~(a) give $r_f\le e_r+2Ke_c$ for a cap, and an auxiliary
face has $r_f=0$. With at most $A+5E$ faces and $2E\le KA+n$,
\[
\hs{(\varphi(w)-I)\xi}\ \le\ A(e_r+2Ke_c)+2K'(A+5E)e_c\ \le\ Ae_r+\bigl[(2K+2K')A+5K'(KA+n)\bigr]e_c .
\]
In a model at precision $e$ both $e_r$ and $e_c$ are at most $e$.
\end{proof}

\begin{remark}\label{rem:linear-filling}
The proof uses the state only through \cref{lem:bounded-transport}. With \cref{lem:transfer-parties}(b) in its place the same
bound holds for $\|(U(w)-I)\Psi\|$ on a purification of the original mixed state, with $e_c$ the actual mirror error of the signed
generators; neither form conditions on a sector or replaces the state. When every mirror is exact, for instance on a maximally
entangled state with transposed mirrors, $e_c=0$ and the bound is the tracial estimate $Ae_r$. At a common precision $e$ the
separately based conjugators of \cref{lem:conjugators} cost $A(A+n)e$, while the tree of faces costs $C_K(A+n)e$. No signed area
and no cancellation between cells is used. For the five Houghton relators $K=12$ and $C_{12}=1633$.
\end{remark}

\begin{fillingproposition}[Far-commutator area]\label{prop:point-fill-sharp}
In the five Johnson relators displayed in \cref{sec:point-source}, put $\sigma_i=a^isa^{-i}$, and let $A^*(R)$ be the largest area of
$[\sigma_i^{\pm1},\sigma_j^{\pm1}]$ with $0\le i,j\le R$ and $|i-j|\ge2$. There is a constant $C$ such that, for every integer $R\ge2$,
\[
 A^*(R)\le CR^{\log_2\lambda},\qquad
 Q(\lambda)=0,\quad Q(t)=t^5-89t^4+824t^3-3416t^2-6672t+19648,
\]
where $\lambda\in(79,80)$ is the largest positive root of $Q$. An admissible area exponent is
\[
 p=\log_2\lambda=6.3064682540119605963\ldots .
\]
\end{fillingproposition}
\begin{proof}
\Cref{app:point-filling} gives the complete six-type filling construction, with every auxiliary relator expanded into the literal five-relator presentation. The recurrence follows the filling method of~\cite{DM-H3}; the explicit matrix and certificates in the appendix establish the exponent stated here.
\end{proof}

\subsection{The window}\label{sec:point-window}

\begin{theorem}[The point window]\label{thm:point-window}
There are constants $c,C>0$ such that, for all sufficiently small $\eps$,
\[
c\,\eps^{-\beta}\ \le\ d_{G_{\mathrm{pt}}}^\to(\eps)\ \le\ C\eps^{-1/4},\qquad
\beta=\frac1{2\log_2\lambda+1}=0.0734595360384268649\ldots .
\]
Here $\lambda$ is the root in \cref{prop:point-fill-sharp}. In particular, every original finite-dimensional mixed state of a strategy of loss at most $\eps$ has $\ED^{A\to B}(\rho)\ge c\eps^{-\beta}$.
By \cref{thm:factor-three} the same window holds for $d^\leftarrow_{G_{\mathrm{pt}}}$. A matching lower exponent $\frac14$ is open
(\cref{sec:open}).
\end{theorem}

The lower bound uses only a fixed visibility certificate, the polynomial far-commutator areas of Houghton's group, the linear
filling charge of \cref{thm:linear-filling} and the $S_3$ extraction of \cref{sec:aggregate}. The upper bound is an explicit
finite model.

\begin{lemma}[A fixed visibility certificate]\label{lem:point-visibility}
Put $\gamma=s\,(asa^{-1})$, a three-cycle. There is a constant $C_0$ such that every model of the point source at precision $e$
has $\hs{(J-I)\xi}\le C_0(\hs{(\gamma-I)\xi}+e)$. Consequently $\hs{(\gamma-I)\xi}\ge\Delta_0>0$ for all sufficiently small $e$.
\end{lemma}

\begin{proof}
With $k=sa$, $l=sb$ and $t=asa^{-1}$, the identity $t=[k,l]$ holds in $H_3$: both sides exchange the second and third points of
the first ray and fix all others, which can be checked on the finitely many points the words move; by Johnson's theorem it has
a fixed finite derivation from the five Houghton relators. In the quotient by $\gamma=1$: $x,z$ commute with $k,l$, hence with $t$,
hence with $s=t^{-1}$, hence with $a=s^{-1}k$; then $[x,aza^{-1}]=1$ becomes $[x,z]=1$, and $[x,z]=J$ gives $J=1$. So $J$ lies in the
normal closure of $\gamma$ and the relators, with a fixed derivation. Evaluating it with \cref{lem:calculus}, every occurrence of
$\gamma$ costs its residual and every fixed conjugator a fixed multiple of $e$. Since $\hs{(J+I)\xi}\le e$, $\hs{(J-I)\xi}\ge2-e$.
\end{proof}

\begin{proof}[Proof of the lower bound]
By \cref{prop:interface}, applied to the universal compiler, every strategy of loss $\eps$ gives Bob unitaries for the six
generators and Alice opposite mirrors, with relator residuals, mirrors, inverse mirrors and state commutators at most
$e=C\sqrt\eps$, on the original state. Round $s$ once to a reflection on each party by $\sgn\Re$; by \cref{lem:rounding} and
\cref{lem:calculus}(e), applied to the fixed relator words, the rounded tuple still has relator residuals, mirrors and state
commutators $O(e)$. This changes observables, not the state. Put $p=\log_2\lambda$ and, for $0\le i<N$,
\[
S_i=a^{3i}sa^{-3i},\qquad T_i=a^{3i+1}sa^{-3i-1},
\]
reflections on each party (Alice using the reversed mirror words). Then $(S_iT_i)^3$ is a conjugate of $(sasa^{-1})^3$ and
$S_iT_i$ a conjugate of $\gamma$, by words of length $O(N)$, so the cubic residuals and raw mirrors are $O(Ne)$ and, by
\cref{lem:point-visibility}, Bob's raw three-cycle visibility is at least $\Delta_0-O(Ne)$. Every cross-copy commutator is a far
commutator $[\sigma_{p'},\sigma_{q'}]$ with $\sigma_j=a^jsa^{-j}$, $0\le p',q'\le3N+1$ and $|p'-q'|\ge2$, a word of length $O(N)$
whose area in the five Houghton relators is at most $C(3N+1)^p$ by \cref{prop:point-fill-sharp}, in the literal presentation and word orientation. Since $p>1$, \cref{thm:linear-filling} bounds its residual by $C(N^p+N)e\le CN^pe$. The additive and the group
commutator of two reflections differ by a state commutator of a word of length $O(N)$, so Bob's raw cross-commutators are at most
$\eta_0\le CN^pe$. \Cref{lem:S3-round} gives exact local $S_3$ representations with
$\kappa\le9\eta_0+144\mu_0+36e_{\mathrm{cub}}\le CN^pe$ and Alice visibility $\Delta\ge\Delta_0/2$ once $Ne$ is small. Take
$N=\lfloor c\,e^{-2/(2p+1)}\rfloor$ with $c$ small: then $N\kappa^2\le Cc^{2p+1}\le\Delta^2/38400$, while $Ne\to0$ because
$2/(2p+1)<1$. For sufficiently small $e$, the floor also gives $N\ge(c/2)e^{-2/(2p+1)}\ge1$. \Cref{cor:S3} gives $\ED^{A\to B}(\rho_{AB})\ge\Delta^2N/16\ge c'e^{-2/(2p+1)}=c''\eps^{-1/(2p+1)}$, the exponent
being halved because $e$ is of order $\sqrt\eps$.
\end{proof}

Charging each cell its own conjugator (\cref{lem:conjugators}) would give only $\eta_0\le CN^{2p}e$ and the exponent
$1/(4p+1)$, with $p=\log_2\lambda$ from \cref{prop:point-fill-sharp}. The area exponent enters only through $\eta_0$;
\cref{sec:open} records what smaller areas would give. The route itself has a limit: no worst-pair estimate $\eta_0\le CN^ue$ with
$u<2$ holds on all original source models, already at the scale of this proof (\cref{thm:raw-obstruction,prop:window-dilution}), so
it cannot give a lower exponent beyond $\frac15$.

\subsection{A collective finite model}\label{sec:point-collective}

The upper bound needs finite tuples of the literal $22$ relators with small defect and little entanglement. The models of
\cite[Proposition~6.5]{DM-H3} spread the single defect of Houghton's corner models over a torus of clocks by rotating the links in
two-dimensional phase planes; on an exterior-algebra quantization they give source precision $O(M^{-3/2})$ at entropy $O(M)$, and
hence only $d\le C\eps^{-1/3}$. The model below does better by taking the logarithm of a \emph{product} of swaps, once, after
representing it.

\begin{proposition}[Collective models]\label{prop:collective}
For every $M\ge48$ there are unitaries for the $22$-relator source on $\C^{D_M}$, $D_M=M^2\,2^{4M+4}$, with $J=-I$, all relators
except H4 exact, and $\opn{\mathrm{H4}-I}\le\pi/A_B+\pi/A_T\le8\pi/M^2$, where $A_B=M(M-1)/2$ and $A_T=1+(M-1)(M-2)/2$. On the maximally
entangled state $\Phi_{D_M}$, with Alice's transposes, every signed generator mirror is exact, and both directional rates equal
$4M+4+2\log M$.
\end{proposition}

\begin{proof}
\emph{Base data.} We use the permutation data of~\cite[\S6]{DM-H3} in its right-action convention: $L=4M+4$ colour coordinates,
clocks $c\in(\Z/M\Z)^2$, the link permutations $A_c,B_c$ of the corner models with the two preliminary exact moves
$\eta_1=(1\ M{+}4)$ at $(M-1,M-1)$ and $\eta_2=(5\ M{+}5)$ at $(0,M-2)$, and $\tau=(1\ 2)$. They satisfy $A_c\tau A_c^{-1}=(2\ 3)$,
$A_c(2\ 3)A_c^{-1}=(3\ 4)$ and $B_c\tau B_c^{-1}=(2\ 3)$ at every clock, and their H4 holonomy is the identity except at the two
sources $q_B=(M-1,M-1)$, where it is $\prod_{t=5}^{M+3}(t\ t{+}M)$, and $q_T=(0,M-2)$, where it is $(5\ M{+}5)$~\cite[proofs of
Propositions~6.1 and~6.5]{DM-H3}.

\emph{One common region.} Let $B$ be the smallest bulk region $R_5$ of~\cite[proof of Proposition~6.5]{DM-H3} and $T$ the region with
source $q_T$, with their labels: for every $t=5,\dots,M+3$ the label $S_t(v)$ of $R_t$ restricted to the vertices of $B$, and the
label $S_T(v)$. Restriction preserves the four regional properties of~\cite[Lemma~6.4]{DM-H3}, because $B\subseteq R_t$ for every $t$: every
base link outside the sources carries each $S_t(v)$, hence their product, to the next label; the regions are connected through
interior edges with zero boundary potentials; base labels avoid $\{1,2\}$ and labels at initial vertices of interior edges avoid
$\{2,3,4\}$; and labels of different regions are disjoint. Let $P_B(v)=\prod_tS_t(v)$, as a product of disjoint transpositions, and
$P_T(v)=S_T(v)$; at the sources these are the two holonomies.

\emph{Represent first, then take one logarithm.} Let $\Pi$ be the qubit-permutation representation on $(\C^2)^{\otimes L}$, and put
$V_R(v)=\Pi(P_R(v))$ and $Q_R(v)=\frac12(I-V_R(v))$ for $R\in\{B,T\}$. These $Q_R$ are orthogonal projections with
$\exp(i\pi Q_R)=V_R$ exactly, covariant along base links ($\Pi(A)Q_R(w)\Pi(A)^*=Q_R(v)$ for an edge from $v$ to $w$), commuting with
one another where both occur, and acting trivially on qubits $2,3,4$. Solve the scalar curl equations
$\mathrm{curl}\,\varphi_R=\pi/A_R-\pi\,1_{q_R}$ on $R$ with zero boundary potentials by a spanning tree of the dual graph, put
$U_c=\prod_R\exp(i\varphi_R(\text{$a$-edge at }c)Q_R(c))$ and $V_c=\prod_R\exp(i\varphi_R(\text{$b$-edge at }c)Q_R(c))$, and define on
$\C^{M^2}\otimes(\C^2)^{\otimes L}$ the operators $a|c,v\rangle=|c-e_1,U_{c-e_1}\Pi(A_{c-e_1})v\rangle$,
$b|c,v\rangle=|c-e_2,V_{c-e_2}\Pi(B_{c-e_2})v\rangle$, $s|c,v\rangle=|c,\Pi(\tau)v\rangle$, $x=X_1$, $z=Z_1$ and $J=-I$.

\emph{The relators.} Since the rotations commute with $\Pi((2\ 3))$ and $\Pi((3\ 4))$, \cite[Lemma~6.3]{DM-H3} gives the four
Houghton relators other than H4 exactly. The new link factors act trivially on qubits $2,3,4$ and the base permutations have the
same first inverse images as in~\cite{DM-H3}, so in the left-operator convention $ax_1a^*=x_2$, $az_1a^*=z_2$, $bx_1b^*=x_2$,
$bz_1b^*=z_2$, $sx_2s^*=x_1$ and $sz_2s^*=z_1$; hence $sa$ and $sb$ fix both lamps and $[u,ava^{-1}]=[u_1,v_2]=1$. Squares, centrality
of the scalar $J$ and $[x,z]=-I=J$ are exact. For H4, the computation of~\cite[Lemma~6.4]{DM-H3} uses only covariance,
commutation, protected supports and the endpoint identity $\exp(i\pi Q)=V$ of its generators, all of which hold here, so the
holonomy at every plaquette is $\exp(i1_B(q)\pi Q_B(q)/A_B)\exp(i1_T(q)\pi Q_T(q)/A_T)$, of operator distance at most
$\pi1_B/A_B+\pi1_T/A_T$ from $I$. Finally $A_B,A_T\ge M^2/4$ for $M\ge4$.

\emph{The state.} On $\Phi_{D_M}$ with Alice's transposes, $(U^T\otimes I)\Phi=(I\otimes U)\Phi$, so every mirror is exact and the
state is pure with entropy $\log D_M$.
\end{proof}

\begin{proof}[Proof of the upper bound in \cref{thm:point-window}]
With the flat state, the collective model is a source model at precision $e_M\le8\pi/M^2$. The upper map of the universal compiler
gives a model of the constraint system at precision $Ce_M$ with the state $I_{12}/12\otimes\rho$, and the boolean return gives a
strategy of loss at most $K_3C^2e_M^2=O(M^{-4})$ on a maximally entangled state of dimension $12D_M$, whose rates are
$4M+4+2\log M+\log12$. Taking $M$ of order $\eps^{-1/4}$ gives $d^\to\le C\eps^{-1/4}$.
\end{proof}

\begin{holonomyproposition}[Fixed-holonomy subdivision from the magnetic-cycle estimate]\label{prop:fixed-holonomy}
Let $A\ge1$, $F_A\cdots F_1=V=e^{iK}$ be finite-dimensional unitaries with $K=K^*$ and $\operatorname{spec}K\subset[-\pi,\pi]$, and let $\rho$ be a density matrix. If $[\rho,V]=0$, then
\[
 \frac1A\sum_j\|(F_j-I)\sqrt\rho\|_2^2\ge4\Tr\bigl(\rho\sin^2(K/(2A))\bigr),
\]
with equality for $F_j=e^{iK/A}$. Without commutation,
\[
 \frac1A\sum_j\|(F_j-I)\sqrt\rho\|_2^2\ge A^{-2}\|(V-I)\sqrt\rho\|_2^2.
\]
For normalized trace and $V=I-2Q$, the infimum is exactly $4\tau(Q)\sin^2(\pi/(2A))$. For the fixed bulk holonomy and area of \cref{prop:collective}, this is $\Theta(M^{-4})$, so changing only its subdivision cannot improve its root-defect order.
\end{holonomyproposition}
\begin{proof}
The complete proof is in \cref{app:fixed-holonomy}. The scalar input is the magnetic-cycle spectrum~\cite[Lemma~2.9]{CTT11} and its equal-distribution minimization~\cite[Proposition~4.6(b)]{BGKLM20}; the arbitrary-state estimate is elementary telescoping. The application fixes the holonomy, representation, visibility and factor-count relation, and does not exclude changing them.
\end{proof}


\begin{remark}\label{rem:point-collective}
The actual tracial H4 error of the collective model is $\Theta(M^{-4})$ in squared normalized Hilbert--Schmidt norm, so the root
precision is genuinely $\Theta(M^{-2})$. The difference from the exterior-algebra models is specific: those take a sum of the
quadratic logarithms of $O(M)$ commuting swaps, whose tracial fluctuations are of size $\sqrt M$, whereas here the logarithm of
their product is one bounded projection. The new links do not factor as a native link times a control, so the native class law
below does not apply to them. An exterior-algebra specialization of the same construction also works, at dimension
$M^22^{8M+8}$.
\end{remark}


\subsection{A conditional law on the calibrated native class}\label{sec:point-native}

The exterior-algebra models of~\cite[Proposition~6.5]{DM-H3}, with $k=2$ flavours and with the Pauli lamps
$x=i\gamma_{1,1}\gamma_{1,2}$, $z=i\gamma_{1,2}\gamma_{1,3}$ built from the Majorana operators of the first site and $J=-I$, satisfy the
$22$ relators with only H4 inexact, $\|\mathrm{H4}-I\|^2_{2,D}\le14.26M^{-3}$, at dimension $M^22^{8M+8}$. Call such a model, for one
$M\ge48$, a \emph{native model} $\mathcal N_M$. On the class of Bob tuples equal to one native model, the demand has a matching
law, but the proof of that law rests on inputs that have not been supplied self-containedly.

\begin{theorem}[Native class law; conditional]\label{thm:native}
Restrict Bob's tuple to one native model $\mathcal N_M$, $M\ge48$, up to a common unitary gauge and arbitrary finite spectators,
and let the original $ABE$ state be arbitrarily correlated. Require actual Alice signed-generator and inverse mirrors at most $e$,
all $22$ actual vector relator residuals at most $e$, and the negative centre; do not assume a flat density, stationarity, clock
diagonality, a pure input, flavour factorization or an independent environment. \emph{Conditional on the inputs listed
below}, the minimum original $\ED^{A\to B}$ on this class is $\Theta(e^{-2/3})$, and on the game $G_{\mathrm{pt}}$, restricted to
the calibrated native class (strategies whose actual native source contract is at most a fixed $C\sqrt\eps$), the minimum demand
is $\Theta(\eps^{-1/3})$. This is not the unrestricted optimum: the collective models of \cref{prop:collective} lie outside the
class, and the unrestricted upper bound $\eps^{-1/4}$ of \cref{thm:point-window} is below this class law.
\end{theorem}

\emph{The conditional inputs.} The lower bound is the aggregate extraction of \cref{thm:aggregate} applied to transported point
lamps, $b_N=9E_B/N+132M_2$, together with four native estimates: the coercivity $M^2e^2+M^3\|(\mathrm{H4}-I)\Psi\|^2\ge c$; a sparse
bad-branch probability $p_{\mathrm{bad}}\le2N/M+4Ne$; the aggregate bound $E_B\le768\,Np_{\mathrm{bad}}$; and hence
$b_N\le13824N/M+27648Ne+396N^3e^2$, which with $N=c\min\{M,e^{-2/3}\}$ gives $\ED\ge c'e^{-2/3}$. The coercivity and the bad-branch
estimates rest on the appendix ``Sharp Houghton cost with moderate flavour growth'' of the authors' manuscript on
immediate delivery~\cite{DM-ID}, whose proofs are not reproduced here, and the bound on $E_B$ has a written derivation that is not included here; the theorem becomes unconditional when self-contained proofs of these are supplied. The aggregate theorem itself
(\cref{thm:aggregate}) is proved. The upper bound uses the literal free-unitary lift $I_{12}\otimes\mathcal N_M$ in the compiled game,
not an assignment correct only modulo the source relations: the recovery polynomial evaluates to $I_{12}\otimes s$ on every
unitary source tuple. The source witness has flat entropy $8M+8+2\log M$; with the amplified flat state the compiled game
witness is maximally entangled of dimension $12M^2\,2^{8M+8}$, belongs to the native class, and has $e\asymp M^{-3/2}$ and
$\ED=8M+8+2\log M+\log12$.

\begin{theorem}[Enlarged native class; conditional]\label{thm:native-enlarged}
Keep exact native clocks and the first two protected flavours, but allow clock-dependent links $L^0_a(c)\otimes U_a(c)$ and
$L^0_b(c)\otimes U_b(c)$ and a clock-diagonal $s_0\otimes S(c)$, with lamps $x_0,z_0\otimes I$, negative $J$, finite noncentral
controls, arbitrary correlations with Alice and the environment, common gauges, and finite coherent direct sums of different
$M$. Require all full actual relator residuals, control defects included, and all actual mirrors, at most $e$. Links must factor
between protected and control factors as stated; arbitrary entangling link operators are not included. Conditional on the same
inputs as \cref{thm:native}, the class law is again $\Theta(e^{-2/3})$, and $\Theta(\eps^{-1/3})$ only for its calibrated
native game restriction.
\end{theorem}

The finite local identity behind the enlarged class is checked exactly: for the generator $H$ of a native link, averaged over the
local Pauli group of the two protected flavours, $\mathbb E_P(H^2)-\mathbb E_P(H)^2=I/8+H_{(+,-)}^2\ge I/8$ on every multiplicity
sector. Its sector estimates give $\sum_Mp_MM^{-3}\le Ce^2$, hence $\sum_Mp_MM^{-1}\le Ce^{2/3}$ by Jensen's inequality, and after
clipping small sectors $b_N\le C_1Ne^{2/3}+C_2Ne+C_3N^3e^2$. The conditional status concerns the operational law, not this
identity.
